% Chapter 4 exercises. Source PDF pages 138--140 / printed pages 125--127.
\MathBookSourcePage{138}
\subsection*{习题}
\label{sec:ch04-exercises}
\setcounter{exerciseitem}{0}

\begin{MBExercise}
\label{exr:ch04-01}
证明对 $|\alpha|\leq m-1$,
\[
\MathBookFormulaID{formula-ch04-exercise-01-derivative-identity}
D_x^\alpha T_y^m u(x)=T_y^{m-|\alpha|}D_x^\alpha u(x)
\qquad\forall u\in C^{|\alpha|}(B).
\]
参见命题~\ref{prop:ch04-averaged-taylor-commutes-with-derivatives} 的证明.
\end{MBExercise}

\begin{MBExercise}
\label{exr:ch04-02}
证明式~\eqref{eq:ch04-taylor-theorem-integral-form} 给出的 Taylor 定理形式.
\end{MBExercise}

\begin{MBExercise}
\label{exr:ch04-03}
验证式~\eqref{eq:ch04-scaled-sobolev-seminorm} 和式~\eqref{eq:ch04-scaled-averaged-taylor-polynomial} 中的断言. 提示: 证明 $T_{d\widehat y}^m u(d\widehat x)=\widehat T_{\widehat y}^m\widehat u(\widehat x)$, 其中
\[
\MathBookFormulaID{formula-ch04-exercise-03-scaled-taylor-polynomial}
\widehat T_{\widehat y}^m\widehat u(\widehat x)
=\sum_{|\alpha|<m}\frac1{\alpha!}D_{\widehat y}^\alpha\widehat u(\widehat y)
(\widehat x-\widehat y)^\alpha.
\]
\end{MBExercise}

\begin{MBExercise}
\label{exr:ch04-04}
证明式~\eqref{eq:ch04-friedrichs-inequality} 等价于
\[
\MathBookFormulaID{formula-ch04-exercise-04-equivalent-friedrichs}
\|u\|_{L^p(\Omega)}
\leq C\left(\left|\int_\Omega u\,dx\right|+|u|_{W_p^1(\Omega)}\right)
\qquad\forall u\in W_p^1(\Omega),
\]
其中常数 $C$ 只依赖于维数 $n$ 和 $\Omega$ 的粗壮度参数 $\gamma$. 提示: 参见式~\MBUnresolvedReference{eq:ch05-03-03}{(5.3.3)} 的证明.
\end{MBExercise}

\begin{MBExercise}
\label{exr:ch04-05}
利用齐次性论证完成定理~\ref{thm:ch04-local-interpolation-error} 一般情形的证明.
\end{MBExercise}

\begin{MBExercise}
\label{exr:ch04-06}
证明三角剖分族 $\{\mathcal T^h\}$ 是拟一致的, 当且仅当它非退化, 且存在与 $h$ 无关的正常数 $c,C$, 使对任意 $K_1,K_2\in\mathcal T^h$,
\[
\MathBookFormulaID{formula-ch04-exercise-06-element-diameter-comparability}
c\operatorname{diam}K_1\leq\operatorname{diam}K_2
\leq C\operatorname{diam}K_1.
\]
\end{MBExercise}

\begin{MBExercise}
\label{exr:ch04-07}
证明不等式~\eqref{eq:ch04-global-interpolation-linfty-error}.
\end{MBExercise}

\MathBookSourcePage{139}
\begin{MBExercise}
\label{exr:ch04-08}
设 $\{a_m\}_{m=1}^\infty$ 是非负数列. 证明若 $1\leq q\leq p$, 则
\[
\MathBookFormulaID{formula-ch04-exercise-08-sequence-norm-monotonicity}
\left(\sum_{m=1}^\infty a_m^p\right)^{1/p}
\leq\left(\sum_{m=1}^\infty a_m^q\right)^{1/q}.
\]
\end{MBExercise}

\begin{MBExercise}
\label{exr:ch04-09}
设 $\{a_m\}_{m=1}^M$ 是有限非负数列. 证明若 $p<q\leq\infty$, 则当 $q<\infty$ 时
\[
\MathBookFormulaID{formula-ch04-exercise-09-finite-sequence-lp-lq}
\left(\sum_{m=1}^Ma_m^p\right)^{1/p}
\leq M^{1/p-1/q}\left(\sum_{m=1}^Ma_m^q\right)^{1/q},
\]
而当 $q=\infty$ 时
\[
\MathBookFormulaID{formula-ch04-exercise-09-finite-sequence-lp-linfty}
\left(\sum_{m=1}^Ma_m^p\right)^{1/p}
\leq M^{1/p}\max_{1\leq m\leq M}a_m.
\]
\end{MBExercise}

\begin{MBExercise}[最小角条件]
\label{exr:ch04-10}
设 $\{\mathcal T^h\}$ 是 $\Omega\subset\mathbb R^2$ 的一族三角剖分. 证明 $\{\mathcal T^h\}$ 非退化, 当且仅当 $\{\mathcal T^h\}$ 中所有三角形的各个角都以某个正常数为下界.
\end{MBExercise}

\begin{MBExercise}
\label{exr:ch04-11}
令 $\widehat K$ 为顶点 $(0,0)$、$(1,0)$、$(0,1)$ 构成的参考三角形, $\zeta\in W_p^m(\widehat K)$, 其中 $m\geq2$, $1\leq p\leq\infty$. 证明
\[
\MathBookFormulaID{formula-ch04-exercise-11-anisotropic-interpolation}
\left\|\frac{\partial}{\partial x_j}(\zeta-I\zeta)\right\|_{L^p(\widehat K)}
\leq C\left|\frac{\partial\zeta}{\partial x_j}\right|_{W_p^m(\widehat K)},
\]
其中 $I\zeta$ 是 $P_{m-1}$ Lagrange 有限元中 $\zeta$ 的插值函数. 其他各向异性插值估计见 Apel (1999).
\end{MBExercise}

\begin{MBExercise}[最大角条件]
\label{exr:ch04-12}
设 $\mathcal T$ 是多边形区域 $\Omega\subset\mathbb R^2$ 的三角剖分, $V_h=\{v\in C^0(\overline\Omega):v|_T\text{ 对所有 }T\in\mathcal T\text{ 都是线性的}\}$. 设 $u\in H^2(\Omega)$, $Iu\in V_h$ 是 $u$ 的线性插值函数. 证明
\[
\MathBookFormulaID{formula-ch04-exercise-12-maximum-angle-estimate}
\|u-Iu\|_{L^2(\Omega)}+h|u-Iu|_{H^1(\Omega)}
\leq\zeta(\theta)h^2|u|_{H^2(\Omega)},
\]
其中 $h=\max_{T\in\mathcal T}\operatorname{diam}T$, $\theta$ 是 $\mathcal T$ 中的最大角, $\zeta$ 是定义在 $[\pi/3,\pi)$ 上的递增正函数. 提示: 使用练习~\ref{exr:ch04-10}.
\end{MBExercise}

\begin{MBExercise}
\label{exr:ch04-13}
对三维 Serendipity 单元导出与第~\ref{sec:ch04-tensor-product-polynomial-approximation} 节中估计类似的估计.
\end{MBExercise}

\begin{MBExercise}
\label{exr:ch04-14}
证明如下 Sobolev 不等式:
\[
\MathBookFormulaID{formula-ch04-exercise-14-sobolev-inequality}
\|v\|_{L^q(\Omega)}\leq\|v\|_{W_p^m(\Omega)}
\qquad\text{当 }m+\frac nq\geq\frac np,
\]
其中要求 $q<\infty$. 提示: 证明 $m$ 阶 Riesz 位势在 $m+n/q\geq n/p$ 时把 $L^p(\Omega)$ 映到 $L^q(\Omega)$, 参见引理~\ref{lem:ch04-sobolev-inequality} 的证明以及 Sobolev (1963 \& 1991).
\end{MBExercise}

\MathBookSourcePage{140}
\begin{MBExercise}
\label{exr:ch04-15}
设 $P$ 是有限维多项式空间, $K$ 是任意包含半径为 $\rho\operatorname{diam}K$ 的球的开集, 其中 $\rho>0$. 证明式~\eqref{eq:ch04-local-inverse-estimate} 成立, 且 $C$ 与 $K$ 无关. 提示: 取球 $B_{\rho h}\subset K\subset B_h$, 并将式~\eqref{eq:ch04-normalized-norm-equivalence} 替换为
\[
\MathBookFormulaID{formula-ch04-exercise-15-ball-norm-equivalence}
\|v\|_{W_p^l(B_h)}\leq C\|v\|_{L^q(B_{\rho h})}
\qquad\forall v\in P.
\]
\end{MBExercise}

\begin{MBExercise}
\label{exr:ch04-16}
证明在二维或更高维中, 非退化三角剖分族 $\{\mathcal T^h\}$ 局部拟一致. 提示: 相邻单元都可以通过一列具有公共面的单元彼此连接.
\end{MBExercise}

\begin{MBExercise}
\label{exr:ch04-17}
考虑第~\ref{sec:ch00-relation-to-difference-methods} 节中一维拟一致网格上的分片线性函数刚度矩阵 $K$. 证明 $K$ 的条件数(Isaacson \& Keller 1966)以 $O(h^{-2})$ 为界. 提示: 使用逆估计.
\end{MBExercise}

\begin{MBExercise}
\label{exr:ch04-18}
完成命题~\ref{prop:ch04-remainder-linfty-bound} 证明中的稠密性论证. 提示: 按构造以及推论~\ref{cor:ch04-averaged-taylor-linfty-bound}, $R^m$ 是从 $W_p^m$ 到 $L^p$ 的连续算子. 对 $W_p^m$ 中的光滑函数 Cauchy 序列, $R^m$ 作用后的序列在 $L^\infty$ 和 $L^p$ 中都是 Cauchy 序列. 证明两个极限相同.
\end{MBExercise}

\begin{MBExercise}
\label{exr:ch04-19}
证明若 $h\geq d/2$, 则引理~\ref{lem:ch04-discrete-sobolev-inequality} 中的估计显然成立. 提示: 应用逆估计~\eqref{eq:ch04-local-inverse-estimate}.
\end{MBExercise}

\begin{MBExercise}
\label{exr:ch04-20}
证明若 $\Omega$ 满足线段条件, 则引理~\ref{lem:ch04-sobolev-inequality} 中的函数 $u$ 在 $\overline\Omega$ 上连续, 且
\[
\MathBookFormulaID{formula-ch04-exercise-20-sobolev-continuity}
\|u\|_{L^\infty(\Omega)}\leq C_{m,n,\gamma,d}\|u\|_{W_p^m(\Omega)}.
\]
提示: 使用注~\ref{rem:ch01-boundary-smooth-density} 中所述 $C^\infty(\overline\Omega)$ 的稠密性以及引理~\ref{lem:ch04-sobolev-inequality} 证明中的论证.
\end{MBExercise}

\begin{MBExercise}
\label{exr:ch04-21}
令 $\widehat K$ 为单位正方形, $K$ 为凸四边形. 证明存在微分同胚 $F:\widehat K\to K$, 使 $F$ 的各分量属于 $Q_1$.
\end{MBExercise}

\begin{MBExercise}
\label{exr:ch04-22}
设 $\{\mathcal T^h\}$ 是凸四边形对多边形区域 $\Omega$ 的非退化剖分族, $V_h$ 是与 $\mathcal T^h$ 关联的 $Q_1$ 等参有限元空间(参见第~\ref{sec:ch04-isoparametric-polynomial-approximation} 节和练习~\ref{exr:ch04-20}), $I^h:H^2(\Omega)\to V_h$ 是节点插值算子. 证明
\[
\MathBookFormulaID{formula-ch04-exercise-22-quadrilateral-isoparametric-error}
\|u-I^h u\|_{L^2(\Omega)}+h|u-I^h u|_{H^1(\Omega)}
\leq Ch^2|u|_{H^2(\Omega)}
\qquad\forall u\in H^2(\Omega).
\]
关于四边形单元的更多内容见 Girault \& Raviart (1979) 以及 Arnold, Boffi \& Falk (2000).
\end{MBExercise}
